% critique.tex — f173r, vv. 877–919 : commentaire critique.
% Parties textuelles du commentaire de l'auteur (travaux antérieurs, Harmattan 2023), texte et notes
% conservés ; répartition, harmonisation et passages signalés : repartition.md.
\cc{895}{βαϋζα τοὐμὸν σῶμα βάλλουσα ψόγῳ.}{La correction de la leçon du \italique{Ravennas} universellement admise, βαϋζα en βάϋζε, est inintelligible. (Les dictionnaires se sont d'ailleurs bien gardés de citer notre vers 895 à l'entrée βαΰζω.)
Faisons l'hypothèse que la fin du mot transmis soit correcte. Le suffixe -υζᾰ a été productif de quelques noms de végétaux en grec\nc{ἄλυζα $=$ ἄλυπον ; κνύζα	 $=$ κόνυζα $=$ σκόνυζα, \italique{conyze} ou \italique{herbe aux puces} ou \italique{encensière} (\latin{erigeron viscosum L.}) ; ἡ μώλυζα, \italique{la tête d'ail} ; ἡ ὄρυζα, \italique{le riz}.} ; la \italique{conyze}, κόνυζα convient parfaitement : les minuscules `κ' et `β' sont très semblables sur de nombreux manuscrits\nc{Ex. : dans les \italique{Euménides} d'Eschyle, lire νεβρόν au lieu de νεκρόν v. 246.} -- dont notre \italique{Ravennas}. 
Le sens s'éclaire ainsi : Mnésiloque se compare à une plante -- une fleur dira-t-on en français -- en proie aux jets de reproches de Critylla. C'est là le point : ψόγῳ à la fin du vers est mis παρὰ προσδοκίαν pour λίθῳ. Or le proverbe grec « jeter une pierre à une herbe » (Apostolius, XIV, 72) :\citComm{πόᾳ λίθον ἐκβάλλει· ἐπὶ τῶν τἀναντία τοῖς ἐπιβούλοις ποιούντων·\\
Il jette une pierre à une herbe\nc{L'explication qui suit de Michel Apostolius devrait plutôt nous conduire à traduire ainsi : « Il retire une pierre à l'aide d'une herbe. » Mais la tentative d'Apostolius de trouver une exégèse savante au proverbe est vouée à l'échec -- comme pour tout ce qui touche à la tradition orale. Le rapprochement que fait Apostolius avec un passage du précieux Élien (\italique{De Natura Animalium}, I, 45), pour astucieux qu'il soit, n'est pas convaincant.}.}}
\cc{910}{ἐγὼ δὲ Μενέλαον ὅσα γ᾽ ἐκ τῶν ἰφύων.}{Scène de reconnaissance. Lorsqu'elle avait reconnu son vieux mari face à elle, la véritable Hélène d'Euripide s'était exclamée (v. 564) :\citComm{Ἐγὼ δὲ Μενέλεῴ γε σέ· οὐδ' ἔχω τί φῶ.\\
Nul ne fut plus pareil à Ménélas. Que dire ?\\
(traduction d'Henri Grégoire.)}
Chez Aristophane, à quoi Hélène reconnaît-elle son naufragé de mari ? « Οὐδ' ἔχω τί φῶ » devient dans le manuscrit de Ravenne : ὅσα γ᾽ ἐκ τῶν ἀφύων, « d'après ton menu fretin ». 
On lit dans le lexique de la \italique{Souda} (ι 778) : ὅσον ἐκ τῶν ἰφύων, « à tes lavandes\nc{Suzanne Amigues identifiait en fait l'ἴφυον avec la nielle des blés (\latin{Agrostemma githago L.}). Voir son article « Quelques légumes de disette chez Aristophane et Plutarque » paru en 1988 dans le \italique{Journal des Savants}, pp. 157-171.} ». 
(C'est cette dernière version que commente le scholiaste, recopié en marge du manuscrit de Ravenne.)
Qu'il s'agisse de \italique{menu fretin} ou de \italique{lavande}, de telles leçons sont ineptes\nc{Il se trouve néanmoins de nombreux défenseurs de la leçon « ἐκ τῶν ἰφύων », leçon qui fait beaucoup rire de savants hellénistes (?!) – décidément très bon public !} !
Colin Austin a bien résumé la situation\nc{«~Observations critiques sur les \italique{Thesmophories} d'Aristophane~», in \italique{Dodone} 19.2 (1990), p. 27.} :
« L'alpha d'ἀφύη étant bref, cette leçon est métriquement impossible. Quant à la ``lavande'' de la \italique{Souda}, le scholiaste y voit une allusion aux herbes de la mère d'Euripide, ce qui semble bien tiré par les cheveux\nc{C'est pourtant cette leçon « bien tirée par les cheveux » qu'Austin fit imprimer dans son édition des \italique{Thesmophories} ; il faut croire que, ici comme en d'autres occasions, son acolyte Olson fut de mauvais conseil.} ! Comme le note H. Grégoire\nc{Tome V de l'édition d'Euripide aux Belles Lettres (1950), p. 65 n. 1.} « les herbes sont ici plus ridicules encore que les anchois », moins ridicules cependant que les \italique{oignons marins} de Gag***\nc{« Une telle saloperie ?.. ou plutôt un Gag, non ? » C'est ce mot d'ordre que je reprends à Austin, modifiant au passage le nom – \latin{nomen omen} ! – du farceur qui va recevoir sa fessée, conformément à ce que réclame la justice !} (τιφύων, « \italique{your squills}... \anglais{here a metaphor for male genitals} ») : cette saloperie-là\nc{Pour réussir à faire perdre son sang-froid au pourtant toujours modéré Colin Austin, il fallait effectivement que ce fût une belle salauderie… J'aurais volontiers laissé la misérable hypo-foutaise « qui se réfute d'elle-même » de Gag*** dans un oubli bien mérité, malheureusement, il m'a été infligé de la voir exhumée et prise au sérieux dans plusieurs articles récents… Enterrons une bonne fois pour toute -- espérons-le ! -- cette charogne : la saloperie de Gag*** présente à la fois le quadruple exploit \ \ \ 1°) \ de ne pas rendre compte de la leçon `ἀφύων' du témoin principal de la pièce : le manuscrit de Ravenne ; \ \ \ 2°) \ d'élucubrer sur un soi-disant sens obscène de τίφυον qui n'est attesté nulle part ; \ \ \ 3°) \ d'être intrinsèquement absurde : si Ménélas est réputé avoir bien peu d'ardeur (raison du départ de l'insatisfaite madame ?) ; comment se fait-il qu'il soit en pleine \italique{érection} et que ce soit à cela qu'Hélène le reconnaisse ? \ \ \ 4°) \ \italique{last but not least}, d'être anti-comique à souhait : un cocu impuissant n'est drôle \italique{que} parce qu'il est impuissant : représentez-le en pleine érection et vous gâterez à coup sûr la farce ! \ \ \ \
Reconnaissons au moins à Gag*** le mérite (pédagogique pour l'édification du lecteur) d'avoir réussi à concentrer en un seul passage les quatre vices cardinaux du \italique{misthophore} de la recherche, savant homme qui se paye – littéralement ! – de mots. Qu'y a-t-il d'étonnant à ce qu'un individu salarié pour palabrer palabre ? Rien, à priori, mais on leur saurait tout-de-même gré d'occuper leur jactance à autre chose que de salir Aristophane.} se réfute d'elle-même. En fait la solution du problème se trouve au v. 935. Là Euripide est accusé d'être un ἱστιορράφος, \italique{un ravaudeur de voiles}. Fritzsche (p. 363)\nc{Voir aussi G. Lange, \italique{Quaest. in Ar. Thesm., Diss. Gott}. (1891), p. 8.} avait bien vu qu'Aristophane se moque ici de l'accoutrement de Ménélas dans l'\italique{Hélène} d'Euripide: tout au long de cette pièce il est fait mention des haillons, des lambeaux de voile dont le roi de Sparte est pitoyablement affublé, comme un gueux, après son naufrage en Égypte (v. 416 τὰς ἐμὰς δυσχλαινίας, 422 ναὸς ἐκβόλοις ἃ ἀμπίσχομαι, 544 sq. ἄγριος δέ τις / μορφὴν ὅδ᾽ ἐστίν, 554 στολήν γ᾽ ἄμορφον, 790 ὥσπερ πτωχὸς, 1204 ὡς ἐσθῆτι δυσμόρφῳ πρέπει, 1282 τῆς ἀχλαινίας, 1297 ἐσθῆτά τ᾽ ἐξάλλαξον). En particulier au v. 423 il s'écrie: « Mes vêtements d'antan, mes habits somptueux, mes anciennes splendeurs, la vague m'a tout pris », πέπλους δὲ τοὺς πρὶν λαμπρά τ᾽ ἀμφιβλήματα / χλιδάς τε πόντος ἥρπασ᾽, et au v. 1079 « les haillons qui m'affublent d'eux-mêmes parleront, en témoins du naufrage », καὶ μὴν τάδ᾽ ἀμφίβληστρα σώματος ῥάκη / ξυμμαρτυρήσει ναυτικῶν ἐρειπίων. Dans les \italique{Grenouilles} (v. 842) Eschyle a bien raison d'appeler Euripide \italique{un faiseur de mendiants}, \italique{un ravaudeur de guenille}, καὶ πτωχοποιὲ καὶ ῥακιοσυρραπτάδη. Pour en revenir au v. 910 des « Thesmo », on est bien tenté de croire que c'est à ses \italique{hardes} que le pauvre naufragé est reconnu par sa femme. »\\ 
Grégoire\nc{Au tome V de son édition d'Euripide aux Belles Lettres, déjà mentionnée. Mais son hypothèse est de douze ans plus ancienne : il l'avait présentée en anglais dans l'article « Modern Greek in Aristophanes » paru dans le premier numéro de \italique{The Link, A Review of Mediaeval and Modern Greek} (1938), pp. 16-20. Une anecdote qui montre comme le monde grec est petit : Henri Grégoire n'est autre que celui qui reprit la suite du grand Alphonse Willems à la chaire de philologie grecque de l'Université libre de Bruxelles en 1909, lorsque ce dernier dut prendre sa retraite.} proposa de lire : `ἐκ τῶν ἀμφίων' ; τὸ ἄμφιον étant une abréviation de ἀμφίεσμα, le \italique{vêtement}. Cette proposition expliquerait la présence des ἀφύων, par simple chute d'une lettre suivie d'iotacisme. En revanche, elle passe à l'as la leçon `ἰφύων' conservée par le scholiaste du manuscrit de Ravenne et par Suidas (ι 778), témoins tout aussi dignes de confiance.
(Re-)prends ton habit Grégoire ! Mais creusons son idée, qui est la bonne ; il nous suffit de trouver un synonyme de vêtement (avec, de préférence, une connotation péjorative) qui puisse rendre compte des deux leçons\nc{Il y avait une autre possibilité de recherche : trouver un mot dont les multiples sens auraient pu être glosés à la fois par ἴφυον et ἀφύη ; les gloses explicatives ayant fini par prendre la place de la leçon originale (le cas s'est déjà vu). Creusant dans cette voie, Austin \& Olson avaient suggéré une émendation très ingénieuse : φυκίων. Émendation ingénieuse sur le plan paléographique, beaucoup moins sur le plan comique…}. 
Ce bon vieil Hésychius ne m'ayant jamais trahi, c'est encore à lui que j'ai posé la question. Cet affable vieillard n'a pas tardé à répondre à ma requête (λ 171) :\citComm{λαίφια· ῥάκη.\\
Laiphia : hardes.}Hésychius nous confirme l'existence d'un diminutif de λαῖφος, \italique{vêtement en lambeaux}, \italique{tissu}, \italique{voile}, ainsi que son sens.\\
Tentons à présent une explication paléographique aux ἀφύων et autres ἰφύων : si l'on fait la simple hypothèse de la disparition de la première lettre du mot λαιφίων\nc{Erreur courante dans notre manuscrit de Ravenne. Citons pour exemple le vers 834 où le manuscrit comporte « δίδοσθαι \textbf{Τ}ηνίοισι » à la place de « δίδοσθαι \textbf{Στ}ηνίοισι ». Cette erreur est d'autant plus fréquente lorsqu'il y a deux lettres identiques consécutives, ainsi au vers 555 du manuscrit de Ravenne : « ὡς \textbf{τ}λεγγίδας λαβοῦσαι » en place de « ὡς \textbf{στ}λεγγίδας λαβοῦσαι ». 
Une telle faute s'explique aisément si l'on imagine le manuscrit copié en onciales et où les mots n'étaient pas même espacés. Dans le mot « λαιφίων », les deux premières lettres, le `λ' et le `α' sont de formes semblables en onciales. La `chute' du `λ' me paraît donc une hypothèse loin d'être invraisemblable.}, le lambda, le scribe se retrouvait alors devant l'inconnu αιφιων, qu'il lisait indifféremment – iotacisme oblige – αιφυων. En butte à ce mot qui n'existe pas, il pouvait alors raisonnablement écrire, selon l'inspiration du moment, ἀφύων ou ἰφύων.}
